\documentclass[11pt]{article}
\usepackage[margin=1in]{geometry}
\usepackage{booktabs,hyperref,xcolor,amsmath}
\definecolor{good}{rgb}{0,0.5,0}
\definecolor{bad}{rgb}{0.8,0,0}
\definecolor{note}{rgb}{0.2,0.2,0.7}
\newcommand{\good}[1]{\textcolor{good}{#1}}
\newcommand{\bad}[1]{\textcolor{bad}{#1}}
\newcommand{\note}[1]{\textcolor{note}{#1}}
\title{RRFA Experiment Log: March 7, 2026\\
\large Re-reading the Original Paper \& Simplifying the Loss}
\author{Internal Notes}
\date{March 7, 2026}
\begin{document}
\maketitle

% ===================================================================
\section{Motivation: Why Re-read the Paper?}
% ===================================================================

After 6 sweeps and dozens of configurations, our best result (Sweep~5, Run~6)
achieved 33.7\% Fujitsu ASR using \texttt{triplet\_full} loss with
\texttt{injection\_aware} masking, contrastive pairs, LoRA 0--20, 300 steps,
and $\gamma_\mathrm{kl}=0.3$. While a significant improvement ($-51$pp from
baseline), several concerns remained:

\begin{enumerate}
  \item \textbf{Complexity}: 8 loss modes, 5 distance functions, 11 masking
        policies, 6-stage ETL pipeline. Too many degrees of freedom for
        confident attribution.
  \item \textbf{dl2rc distance}: $L_2 + 10 \cdot \mathrm{ReLU}(\cos)$ was
        introduced to fix a zero-gradient problem, but margins (1.2/1.6)
        were tuned for cosine distance $\in [0,2]$, while $L_2$ on 4096-dim
        hidden states can reach 10--100+. Margins would saturate immediately.
  \item \textbf{Prior ``failures'' of paper methods}: We recorded that
        ReLU(cos\_sim) loss ``FAILED'' (0\% Fujitsu improvement) and the
        coefficient schedule ``FAILED'' (washes out rerouting). But had we
        implemented them correctly?
\end{enumerate}

This prompted a careful re-reading of the original Circuit Breakers paper
\cite{zou2024circuit}, Algorithm~1.

% ===================================================================
\section{What the Paper Actually Does}
% ===================================================================

\subsection{Algorithm 1: LoRRA with RR Loss}

The paper's loss is remarkably simple:
\begin{align}
  \mathcal{L}_s &= \mathrm{ReLU}\bigl(\cos\!\operatorname{sim}(\mathrm{rep}_\mathcal{M}(x_s),\; \mathrm{rep}_{\mathcal{M}_\mathrm{cb}}(x_s))\bigr) \label{eq:rr} \\
  \mathcal{L}_r &= \bigl\| \mathrm{rep}_\mathcal{M}(x_r) - \mathrm{rep}_{\mathcal{M}_\mathrm{cb}}(x_r) \bigr\|_2 \label{eq:retain} \\
  \mathcal{L} &= c_s \cdot \mathcal{L}_s + c_r \cdot \mathcal{L}_r \label{eq:total}
\end{align}
with the coefficient schedule:
\begin{align}
  c_s &= \alpha \cdot \frac{t}{2T} \qquad \text{(reroute: } 0 \to \alpha/2\text{)} \label{eq:cs} \\
  c_r &= \alpha \cdot \left(1 - \frac{t}{2T}\right) \qquad \text{(retain: } \alpha \to \alpha/2\text{)} \label{eq:cr}
\end{align}

\noindent Two terms. No margins, no triplet, no cluster centers, no KL divergence.

\subsection{Paper's Exact Hyperparameters (Llama-3-8B-Instruct)}

\begin{center}
\begin{tabular}{ll}
\toprule
Parameter & Value \\
\midrule
LoRA layers & 0--20 (all linear layers) \\
CB target layers & 10, 20 \\
Steps & 150 \\
Batch size & 16 \\
$\alpha$ & 10 \\
Training time & $\sim$20 min (1$\times$ A100) \\
\bottomrule
\end{tabular}
\end{center}

\subsection{Paper's Loss Ablation Results}

\begin{center}
\begin{tabular}{lccc}
\toprule
Loss variant & Converged? & Avg ASR $\downarrow$ & MT-Bench \\
\midrule
RMU ($L_2$ to random vector) & \bad{No} & --- & --- \\
RandC (random centered unit vec) & \bad{No} & --- & --- \\
RandP (random positive unit vec) & Yes & 9.7\% & 8.0 \\
\good{RR (ReLU cos\_sim)} & \good{Yes} & \good{2.5\%} & \good{8.0} \\
\bottomrule
\end{tabular}
\end{center}

RR was the clear winner. The alternatives either failed to converge or
were 4$\times$ worse.

\subsection{Token Selection for CB Loss}

Section~A.1, Design Consideration~2: \emph{``For enhanced robustness, we apply
the circuit-breaking loss to both the user and assistant text within the
circuit breaker set for large language models.''}

The paper applies CB loss to \textbf{all tokens} in the harmful set ---
not just the assistant response, not just the injection span.

% ===================================================================
\section{Critical Discovery: Our Coefficient Schedule Was Backwards}
% ===================================================================

Our \texttt{get\_dual\_coefficients()} implemented:
\begin{align*}
  c_s &: 1 \to 0 \qquad \text{(reroute \emph{decays} to zero)} \\
  c_r &: 0 \to 1 \qquad \text{(retain \emph{ramps up})}
\end{align*}

The paper's schedule is the \textbf{opposite}:
\begin{align*}
  c_s &: 0 \to \alpha/2 \qquad \text{(reroute \emph{ramps up})} \\
  c_r &: \alpha \to \alpha/2 \qquad \text{(retain \emph{ramps down})}
\end{align*}

\textbf{This is why we recorded ``coefficient schedule FAILED --- washes out
rerouting.''} We were literally killing the rerouting signal over the course of
training. The paper does the opposite: it starts with pure retention (stabilize
the model), then gradually introduces rerouting.

At step $T$: $c_s = \alpha/2 = 5$, $c_r = \alpha/2 = 5$. Both terms contribute
equally at the end. The rerouting signal is at its strongest when it matters
most --- after the model has been stabilized by early retention.

% ===================================================================
\section{Re-examining ``ReLU(cos\_sim) Failed''}
% ===================================================================

In Sweep~3/4, we recorded: \emph{``legacy\_cb BROKEN for Fujitsu. 85\% ASR =
baseline. ReLU(cos\_sim) provides $\sim$zero gradient at LoRA scale.''}

\subsection{The Zero-Gradient Problem at LoRA Initialization}

At LoRA initialization ($B=0$), the LoRA model produces identical
representations to the frozen model. The gradient of
$\cos\!\operatorname{sim}(a, b)$ with respect to $b$ when $a = b$ is:
\[
  \frac{\partial}{\partial b} \frac{a \cdot b}{\|a\| \|b\|} \bigg|_{a=b}
  = \frac{a}{\|a\|\|b\|} - \frac{b(a \cdot b)}{\|a\|\|b\|^3} \bigg|_{a=b} = 0
\]

This is a genuine mathematical zero. The gradient of $\mathrm{ReLU}(\cos\!\operatorname{sim})$
at initialization is exactly zero.

\subsection{How the Paper Works Despite This}

The paper's coefficient schedule resolves this: at step 0, $c_s = 0$. The
rerouting loss \emph{is not applied at all} initially. During the early
retention-dominated phase, weight decay and numerical noise create infinitesimal
perturbations that break the exact $a = b$ symmetry. Once
$\cos\!\operatorname{sim}$ drops from 1.0000 to 0.9999, the gradient becomes
non-zero and training takes off.

\textbf{Our Sweep~3/4 used the backwards schedule} ($c_s = 10$ at step 0),
which tried to push rerouting at maximum strength when the gradient was exactly
zero. This explains the observation: \texttt{cos\_sim=1.0000} and
\texttt{grad\_norm=0.0007} at \emph{every} checkpoint.

\subsection{The Paper DOES Use LoRA}

A key misconception from our earlier analysis: we wrote \emph{``Original CB
paper uses full fine-tuning.''} This is incorrect. Algorithm~1 is explicitly
titled \emph{``LoRRA (RepE method) with Representation Rerouting (RR) Loss.''}
Section~A.1: \emph{``To ensure greater stability and improved retention
performance, we employ LoRA tuning.''}

The paper uses the \textbf{same setup we have}: LoRA layers 0--20, CB layers
10+20, Llama-3-8B-Instruct. There is no LoRA/loss incompatibility --- our
schedule was simply wrong.

% ===================================================================
\section{The Complexity Progression (and Why Simplicity Wins)}
% ===================================================================

Our loss evolved through a chain of workarounds for a misdiagnosed root cause:

\begin{center}
\small
\begin{tabular}{lll}
\toprule
Step & What we did & Why \\
\midrule
1 & Used \texttt{legacy\_cb} (ReLU cos\_sim) & Paper's loss \\
2 & Observed 0\% improvement & Blamed loss/LoRA incompatibility \\
3 & Switched to \texttt{triplet\_full} & Needed non-zero gradients \\
4 & Added cluster centers, cross-terms & Triplet needed anchors \\
5 & Added margins (1.2/1.6) & Triplet needed convergence targets \\
6 & Added KL divergence ($\gamma_\mathrm{kl}$) & Needed benign preservation \\
7 & Added \texttt{dl2rc} distance & cos\_sim still had zero grad \\
8 & Realized margins broken for dl2rc & $L_2 \gg 2$, margins saturate \\
\bottomrule
\end{tabular}
\end{center}

\textbf{Steps 2--8 were all consequences of a wrong coefficient schedule.}
With the correct schedule (reroute ramps UP), the paper's 2-term loss should
work out of the box. Every additional component we added was solving a
downstream symptom, not the root cause.

This is a textbook example of \emph{accidental complexity}: each fix was
locally rational given the evidence, but the chain of fixes compounded into
a system with $>$10 hyperparameters, any of which could silently break training.

% ===================================================================
\section{Implementation: \texttt{original\_rr} Mode}
% ===================================================================

\subsection{Core Loss}

Added as a single mode in \texttt{losses.py} using existing primitives
(\texttt{reroute\_loss\_relu\_cos} + \texttt{retain\_loss\_l2}):

\begin{equation}
  \mathcal{L} = c_s \cdot \underbrace{\mathrm{ReLU}(\cos\!\operatorname{sim}(\mathrm{rep}_\mathrm{frozen}^h, \mathrm{rep}_\mathrm{cb}^h))}_{\text{reroute: push harmful orthogonal}}
  + c_r \cdot \underbrace{\|\mathrm{rep}_\mathrm{frozen}^b - \mathrm{rep}_\mathrm{cb}^b\|_2}_{\text{retain: keep benign close}}
  + \gamma_\mathrm{kl} \cdot \underbrace{D_\mathrm{KL}(p_\mathrm{cb}^b \| p_\mathrm{frozen}^b)}_{\text{optional: preserve logits}}
\end{equation}

\begin{itemize}
  \item $\gamma_\mathrm{kl} = 0$ (default): paper exact, 2 terms
  \item $\gamma_\mathrm{kl} > 0$: paper + output-level preservation, 3 terms
\end{itemize}

\subsection{Masking}

\begin{center}
\begin{tabular}{lll}
\toprule
Component & Harmful mask & Benign mask \\
\midrule
Reroute ($\mathcal{L}_s$) & \texttt{loss\_mask} (injection\_aware) & --- \\
Retain ($\mathcal{L}_r$) & --- & \texttt{attention\_mask} (all tokens) \\
KL ($D_\mathrm{KL}$) & --- & \texttt{attention\_mask}, \textbf{NOT loss\_mask} \\
\bottomrule
\end{tabular}
\end{center}

The paper applies CB to all harmful tokens. Our \texttt{injection\_aware} mask
is a refinement (0.5 on injection span, 1.0 on post-injection assistant).
Both approaches are testable via the existing \texttt{--policy-override} flag.

\subsection{What \texttt{original\_rr} Does NOT Use}

\begin{center}
\begin{tabular}{ll}
\toprule
Removed & Reason \\
\midrule
Triplet formulation & Paper uses direct push/pull, not relative distances \\
Cluster centers & No anchors needed --- frozen model IS the anchor \\
Margins (1.2/1.6) & ReLU(cos\_sim) naturally saturates at 0 \\
dl2rc distance & Paper's cos\_sim + L2 retain is sufficient \\
Contrastive pairs & Paper doesn't use them \\
Pooling & Per-token loss, no aggregation needed \\
\bottomrule
\end{tabular}
\end{center}

% ===================================================================
\section{Sweep 7 Plan}
% ===================================================================

\subsection{Runs}

\begin{center}
\begin{tabular}{llcccl}
\toprule
\# & Mode & $\gamma_\mathrm{kl}$ & Steps & $\alpha$ & Purpose \\
\midrule
1 & \texttt{original\_rr} & 0.0 & 150 & 10 & Paper exact baseline \\
2 & \texttt{original\_rr} & 0.3 & 150 & 10 & Paper + KL preservation \\
\bottomrule
\end{tabular}
\end{center}

Both use: LoRA 0--20, CB layers 10+20, \texttt{injection\_aware} harmful mask,
\texttt{attention\_mask} for benign, batch 16 (4$\times$4).

\subsection{What to Watch in Logs}

\begin{itemize}
  \item \textbf{cos\_sim at step 0}: Should be 1.0000 (identical reps). $c_s = 0$, so zero gradient is expected.
  \item \textbf{cos\_sim at step 10--20}: Should start dropping. If still 1.0000 at step 50, something is wrong.
  \item \textbf{L\_rr trajectory}: Should decrease from $\sim$1.0 toward 0. Indicates harmful reps becoming orthogonal.
  \item \textbf{L\_ret trajectory}: Should start near 0 and increase slightly as rerouting creates benign drift.
  \item \textbf{grad\_norm}: Should be non-zero by step 5--10 (weight decay breaks symmetry).
\end{itemize}

\subsection{Key Hypotheses}

\begin{enumerate}
  \item \textbf{H1}: \texttt{original\_rr} with $\gamma_\mathrm{kl}=0$ will match or beat
        our Sweep~5 best (33.7\% ASR) with far fewer hyperparameters.
  \item \textbf{H2}: If the paper's loss works on our data with the correct schedule,
        this confirms the root cause was the backwards schedule, not a loss/LoRA
        incompatibility.
  \item \textbf{H3}: Adding KL ($\gamma_\mathrm{kl}=0.3$) will improve benign
        preservation without hurting harmful ASR.
  \item \textbf{H4}: If \texttt{original\_rr} underperforms, the difference is
        attributable to our agent-specific data (tool calls, injection patterns)
        vs the paper's text safety data, not the loss function.
\end{enumerate}

% ===================================================================
\section{Retrospective: What Each Sweep Taught Us}
% ===================================================================

\begin{center}
\small
\begin{tabular}{clp{8cm}}
\toprule
Sweep & Date & Key Lesson \\
\midrule
1 & Feb 28 & Pooling bug: \texttt{[B,T]} mask broadcasts as \texttt{[B,1,H]} when $T=H=4096$. 99.4\% sparse. \\
2 & Mar 3 & Layer depth matters: 10+20 \good{$\gg$} 20--31 (14pp). Lossmask fix enables benign preservation. \\
3 & Mar 4 & \texttt{legacy\_cb} gradient $\approx 0$ at LoRA scale. Triplet\_full provides non-zero gradient. \\
4 & Mar 4 & Reproduction within $\pm 2$pp. NTP eval ceiling = 21\%. \\
5 & Mar 5 & \texttt{injection\_aware} is the key breakthrough ($-27.5$pp). Contrastive pairs alone don't help. \\
6 & --- & (Planned but superseded by paper re-analysis) \\
7 & Mar 7 & \note{Paper re-read: coefficient schedule was backwards. Simplify to 2-term loss.} \\
\bottomrule
\end{tabular}
\end{center}

% ===================================================================
\section{Discussion: Why the Exploration Was Valuable}
% ===================================================================

Although we ultimately returned to a simpler loss, the intermediate sweeps
produced insights that would not have emerged from the paper's loss alone:

\begin{enumerate}
  \item \textbf{injection\_aware masking} (Sweep~5): The paper applies CB to all
        tokens. Our focused masking (0.5 injection, 1.0 post-injection
        assistant) gave $-27.5$pp improvement in a single change. This is a
        genuine contribution beyond the paper --- the paper's data (harmful text
        completions) doesn't have an ``injection boundary'' to exploit.

  \item \textbf{Benign data filtering bug} (Sweep~7 prep): With
        \texttt{injection\_aware}, benign traces get all-zero \texttt{loss\_mask}.
        A zero-mask filter in data loading silently dropped ALL benign samples.
        Sweep~5's result was achieved with only $\sim$250 resisted AD traces
        as benign data. The true benign set ($\sim$16k traces) was never used.
        Fixing this may further improve results.

  \item \textbf{Per-token vs pooled representations}: Sweeps 1--4 used pooled
        representations with the legacy pooling bug ($T=H$ broadcast).
        Understanding this bug led to correct pooling and ultimately to the
        realization that per-token loss (no pooling at all) is simpler and
        more robust.

  \item \textbf{Failure mode taxonomy}: Three distinct failure modes identified
        (reasoning text, code generation, direct answer). These inform future
        data augmentation, not loss design.

  \item \textbf{Training data imbalance}: The 33 ``hard case'' failures
        concentrate in the minority direction (5.6\% of data, 89\% failure
        rate). This is a data problem, not a loss problem --- further evidence
        that simplifying the loss is the right move.
\end{enumerate}

\subsection{What We'd Keep From the Complex Pipeline}

\begin{itemize}
  \item \texttt{injection\_aware} masking (our key contribution)
  \item Optional KL preservation (cheap insurance against benign degradation)
  \item \texttt{next\_tool\_prediction} eval mode (correct benign metric)
  \item The 6-stage ETL pipeline (necessary for agent trace data, regardless of loss)
\end{itemize}

\subsection{What We'd Drop}

\begin{itemize}
  \item Triplet formulation, cluster centers, cross-terms
  \item dl2rc distance (solving a non-problem under correct schedule)
  \item Margins (ReLU saturates naturally)
  \item 5 of 8 loss modes (\texttt{original\_rr} + \texttt{triplet\_full} as fallback)
  \item Contrastive pairs (no measurable benefit without \texttt{injection\_aware})
\end{itemize}

% ===================================================================
\section{Open Questions}
% ===================================================================

\begin{enumerate}
  \item Does the paper's loss + our \texttt{injection\_aware} mask outperform
        the paper's loss + the paper's all-token mask? (Testable via
        \texttt{--policy-override cb\_full\_sequence})
  \item Does fixing the benign data filtering bug (loading all $\sim$16k benign
        traces instead of $\sim$250) change results significantly?
  \item Can 150 steps (paper default) match 300 steps (our Sweep~5 best) now
        that the schedule is correct?
  \item Does the paper's retain set composition matter? (UltraChat + XSTest vs
        our benign Fujitsu + AgentDojo traces)
\end{enumerate}

\begin{thebibliography}{1}
\bibitem{zou2024circuit}
A.~Zou, L.~Phan, J.~Wang, D.~Duber, M.~Mu, and D.~Hendrycks.
\newblock Improving alignment and robustness with circuit breakers.
\newblock \emph{arXiv preprint arXiv:2406.04313}, 2024.
\end{thebibliography}

\end{document}
